\documentclass[a4paper]{arrowhead}

\usepackage[yyyymmdd]{datetime}
\usepackage{etoolbox}
\usepackage[utf8]{inputenc}
\usepackage{multirow}

\renewcommand{\dateseparator}{-}

\setlength{\parskip}{1em}

%% Special references
\newcommand{\fref}[1]{{\textcolor{ArrowheadBlue}{\hyperref[sec:functions:#1]{#1}}}}
\newcommand{\mref}[1]{{\textcolor{ArrowheadPurple}{\hyperref[sec:model:#1]{#1}}}}
\newcommand{\pdef}[1]{{\textcolor{ArrowheadGrey}{#1\label{sec:model:primitives:#1}\label{sec:model:primitives:#1s}\label{sec:model:primitives:#1es}}}}
\newcommand{\pref}[1]{{\textcolor{ArrowheadGrey}{\hyperref[sec:model:primitives:#1]{#1}}}}

\newrobustcmd\fsubsection[3]{
  \addtocounter{subsection}{1}
  \addcontentsline{toc}{subsection}{\protect\numberline{\thesubsection}interface \textcolor{ArrowheadBlue}{#1}}
  \renewcommand*{\do}[1]{\rref{##1},\ }
  \subsection*{
    \thesubsection\quad
    interface
    \textcolor{ArrowheadBlue}{#1}
    (\notblank{#2}{\mref{#2}}{})
    \notblank{#3}{: \mref{#3}}{}
  }
  \label{sec:functions:#1}
}
\newrobustcmd\msubsection[2]{
  \addtocounter{subsection}{1}
  \addcontentsline{toc}{subsection}{\protect\numberline{\thesubsection}#1 \textcolor{ArrowheadPurple}{#2}}
  \subsection*{\thesubsection\quad#1 \textcolor{ArrowheadPurple}{#2}}
  \label{sec:model:#2} \label{sec:model:#2s} \label{sec:model:#2es}
}
\newrobustcmd\msubsubsection[3]{
  \addtocounter{subsubsection}{1}
  \addcontentsline{toc}{subsubsection}{\protect\numberline{\thesubsubsection}#1 \textcolor{ArrowheadPurple}{#2}}
  \subsubsection*{\thesubsubsection\quad#1 \textcolor{ArrowheadPurple}{#2}}
  \label{sec:model:#2} \label{sec:model:#2s}
}
%%

\begin{document}

%% Arrowhead Document Properties
\ArrowheadTitle{query-multi} % XXX = ServiceName 
\ArrowheadServiceID{query-multi} % ID name of service
\ArrowheadType{Service Description}
\ArrowheadTypeShort{SD}
\ArrowheadVersion{4.4.0} % Arrowhead version X.Y.Z, e..g. 4.4.1
\ArrowheadDate{\today}
\ArrowheadAuthor{Tamás Bordi} % Corresponding author e.g. Jerker Delsing
\ArrowheadStatus{RELEASE} % e..g. RELEASE, RELEASE CONDIDATE, PROTOTYPE
\ArrowheadContact{tbordi@aitia.ai} % Email of corresponding author
\ArrowheadFooter{\href{www.arrowhead.eu}{www.arrowhead.eu}}
\ArrowheadSetup
%%

%% Front Page
\begin{center}
  \vspace*{1cm}
  \huge{\arrowtitle}

  \vspace*{0.2cm}
  \LARGE{\arrowtype}
  \vspace*{1cm}

  %\Large{Service ID: \textit{"\arrowid"}}
  \vspace*{\fill}

  % Front Page Image
  %\includegraphics{figures/TODO}

  \vspace*{1cm}
  \vspace*{\fill}

  % Front Page Abstract
  \begin{abstract}
    This document provides service description for the \textbf{query-multi} service. 
  \end{abstract}

  \vspace*{1cm}

%   \scriptsize
%   \begin{tabularx}{\textwidth}{l X}
%     \raisebox{-0.5\height}{\includegraphics[width=2cm]{figures/artemis_logo}} & {ARTEMIS Innovation Pilot Project: Arrowhead\newline
%     THEME [SP1-JTI-ARTEMIS-2012-AIPP4 SP1-JTI-ARTEMIS-2012-AIPP6]\newline
%     [Production and Energy System Automation Intelligent-Built environment and urban infrastructure for sustainable and friendly cities]}
%   \end{tabularx}
%   \vspace*{-0.2cm}
 \end{center}

\newpage
%%

%% Table of Contents
\tableofcontents
\newpage
%%

\section{Overview}
\label{sec:overview}
This document describes the \textbf{query-multi} service, which enables autonomous service look up for multiple services by one request, therefore it is an integral part of the implementation of service discovery requirements in Service Registry Mandatory Core System. Examples of this interaction is a core system that needs the available information for services that fulfill certain requirements. The look up process itself is exactly the same as in case of the \textit{query} service.

The rest of this document is organized as follows.
In Section \ref{sec:functions}, we describe the abstract message functions provided by the service.
In Section \ref{sec:model}, we end the document by presenting the data types used by the mentioned functions.

\newpage

\subsection{How This Service Is Meant to Be Used}
The given core system is required to submit a list of service query forms, in which the requirements for each systems are defined.
\subsection{Important Delimitations}
\label{sec:delimitations}

The query data must meet the following criteria:

\begin{itemize}
    \item Service definition can contain maximum 63 character of letters (english alphabet), numbers and dash (-), and has to start with a letter (also cannot ends with dash).
    \item Interface names have to follow the \texttt{Protocol-SecurityType-MimeType} format.
    \item Security types could be only \texttt{NOT\_SECURE}, \texttt{CERTIFICATE} or \texttt{TOKEN} .
\end{itemize}

\subsection{Access policy}
\label{sec:accesspolicy}

Available only for the following core systems: \textit{Authorizationt, Orchestrator, Certificate Authority, Choreographer, Device Registry, Event Handler, HawkBit Configuration Manager, MSCV, Plant Description Engine, QoS-Monitor, System Registry}

\newpage

\section{Service Interface}
\label{sec:functions}

This section describes the interfaces to the service. The \textbf{query-multi} service is used to looking for multiple services. A service could contain various metadata as well as a physical endpoint. The various parameters are representing the necessary system and service input information.
In particular, each subsection names an interface, an input type and an output type, in that order.
The input type is named inside parentheses, while the output type is preceded by a colon.
Input and output types are only denoted when accepted or returned, respectively, by the interface in question. All abstract data types named in this section are defined in Section 3.

The following interfaces are available.

\fsubsection{HTTP/TLS/JSON}{ServiceQueryForm}{ServiceQueryResult}

\begin{table}[ht!]
  \centering
  \begin{tabular}{|l|l|l|l|}
    \rowcolor{gray!33} Profile ype & Type & Version \\ \hline
    Transfer protocol & HTTP & 1.1 \\ \hline
    Data encryption & TLS & 1.3 \\ \hline
    Encoding & JSON & RFC 8259 \cite{rfc8259} \\ \hline
    Compression & N/A & - \\ \hline
  \end{tabular}
  \caption{HTTP/TLS/JSON communication details.}
  \label{tab:comunication_semantics_profile}
\end{table}

\clearpage

\section{Information Model}
\label{sec:model}

Here, all data objects that can be part of the \textbf{query-multi} service
provides to the hosting System are listed in alphabetic order.
Note that each subsection, which describes one type of object, begins with the \textit{struct} keyword, which is used to denote a collection of named fields, each with its own data type.
As a complement to the explicitly defined types in this section, there is also a list of implicit primitive types in Section \ref{sec:model:primitives}, which are used to represent things like hashes and identifiers.

\msubsection{struct}{ServiceQueryFormList}

\begin{table}[ht!]
\begin{tabularx}{\textwidth}{| p{5cm} | p{3cm} | p{2cm} | X |} \hline
\rowcolor{gray!33} Field & Type & Mandatory & Description \\ \hline
forms &\pref{Array}$<$\pref{Object}$>$ & yes & List of service requirements form. \\ \hline
\end{tabularx}
\end{table}

\msubsubsection{struct}{forms}

\begin{table}[ht!]
\begin{tabularx}{\textwidth}{| p{5cm} | p{3cm} | p{2cm} | X |} \hline
\rowcolor{gray!33} Field & Type & Mandatory & Description \\ \hline
interfaceRequirements &\pref{Array}$<$\pref{Interface}$>$ & no & Names of the required interfaces. If specified at least one of the interfaces must match for having result(s). \\ \hline
maxVersionRequirement &\pref{Number} & no & Required maximum version of the service. If specified version must be equals or lower or having result(s). Ignored if \textit{versionRequirement} is specified. \\ \hline
metadataRequirements &\pref{Metadata} & no & Service metadata requirements. If specified the whole content of
the map must match for having result(s) \\ \hline
minVersionRequirement &\pref{Number} & no & Required minimum version of the service. If specified version must be equals or higher or having result(s). Ignored if \textit{versionRequirement} is specified. \\ \hline
pingProviders &\pref{Boolean} & no & Whether or not the provider should be pinged. If \textit{true} only the responding providers will comply. \\ \hline
securityRequirements &\pref{Array}$<$\pref{SecureType}$>$ & no & Types of the required security levels. If specified at least one of the types must match for having result(s). \\ \hline
serviceDefinitionRequirement &\pref{Name} & yes & Identifier of the service. \\ \hline
versionRequirement &\pref{Number} & no & Required version of the service. If specified version must match for having result(s) \\ \hline
\end{tabularx}
\end{table}

\msubsubsection{struct}{Metadata}

A JSON \pref{Object} which maps \pref{String} key-value pairs

\clearpage

\msubsection{struct}{ServiceQueryResultList}

\begin{table}[ht!]
\begin{tabularx}{\textwidth}{| p{4.25cm} | p{3.5cm} | X |} \hline
\rowcolor{gray!33} Field & Type      & Description \\ \hline
results & \pref{Array}$<$\pref{Object}$>$     & List of service query results \\ \hline
\end{tabularx}
\end{table}

\msubsubsection{struct}{results}

\begin{table}[ht!]
\begin{tabularx}{\textwidth}{| p{3cm} | p{6cm} | X |} \hline
\rowcolor{gray!33} Field & Type      & Description \\ \hline
serviceQueryData & \pref{Array}$<$\pref{ServiceRegistryResponse}$>$     & List of service instance results \\ \hline
unfilteredHits & \pref{Number} & Number of service instances not included by \textit{serviceQueryData} due to the specified filters. \\ \hline
\end{tabularx}
\end{table}

\msubsubsection{struct}{ServiceRegistryResponse}

\begin{table}[ht!]
\begin{tabularx}{\textwidth}{| p{4.25cm} | p{3.5cm} | X |} \hline
\rowcolor{gray!33} Field & Type      & Description \\ \hline
createdAt & \pref{DateTime} & Service instance record was created at this UTC timestamp. \\ \hline
endOfValidity & \pref{DateTime} & Service is available until this UTC timestamp. \\ \hline
id & \pref{Number} & Identifier of the service instance \\ \hline
interfaces & \pref{Array}$<$\pref{Object}$>$     & List of interfaces the service supports. \\ \hline
provider & \pref{Object} & Descriptor of the provider system record. \\ \hline
secure &\pref{SecureType}  & Type of security the service uses. \\ \hline
serviceDefinitionResponse & \pref{Object} & Descriptor of the serviceDefinition record. \\ \hline
serviceUri &\pref{URI}         & URI of the service. \\ \hline
metadata & \pref{Metadata}     & Service metadata \\ \hline
updatedAt & \pref{DateTime} & Service instance record was modified at this UTC timestamp. \\ \hline
version &\pref{Version}     & Version of the service. \\ \hline
\end{tabularx}
\end{table}

\msubsubsection{struct}{interfaces}

\begin{table}[ht!]
\begin{tabularx}{\textwidth}{| p{4.25cm} | p{3.5cm} | X |} \hline
\rowcolor{gray!33} Field & Type      & Description \\ \hline
createdAt & \pref{DateTime} & Interface instance record was created at this UTC timestamp. \\ \hline
id & \pref{Number} & Identifier of the interface instance \\ \hline
interfaceName &\pref{Interface}  & Specified name of the interface. \\ \hline
updatedAt & \pref{DateTime} & Interface instance record was modified at this UTC timestamp. \\ \hline
\end{tabularx}
\end{table}

\clearpage

\msubsubsection{struct}{provider}

\begin{table}[ht!]
\begin{tabularx}{\textwidth}{| p{4.25cm} | p{3.5cm} | X |} \hline
\rowcolor{gray!33} Field & Type      & Description \\ \hline
address & \pref{String} & Network address. \\ \hline
authenticationInfo & \pref{String}     & Public key of the client certificate. \\ \hline
createdAt & \pref{DateTime} & System instance record was created at this UTC timestamp. \\ \hline
id & \pref{Number} & Identifier of the system instance \\ \hline
metadata & \pref{Metadata}     & Metadata \\ \hline
port & \pref{PortNumber} & Port of the system. \\ \hline
systemName &\pref{Name}  & Name of the system. \\ \hline
updatedAt & \pref{DateTime} & System instance record was modified at this UTC timestamp. \\ \hline
\end{tabularx}
\end{table}

\msubsubsection{struct}{serviceDefinition}

\begin{table}[ht!]
\begin{tabularx}{\textwidth}{| p{4.25cm} | p{3.5cm} | X |} \hline
\rowcolor{gray!33} Field & Type      & Description \\ \hline
createdAt & \pref{DateTime} & Service definition instance record was created at this UTC timestamp. \\ \hline
id & \pref{Number} & Identifier of the service definition instance \\ \hline
serviceDefinition &\pref{Name}  & Name of the service definition. \\ \hline
updatedAt & \pref{DateTime} & Service definition instance record was modified at this UTC timestamp. \\ \hline
\end{tabularx}
\end{table}

\subsection{Primitives}
\label{sec:model:primitives}

Types and structures mentioned throughout this document that are assumed to be available to implementations of this service.
The concrete interpretations of each of these types and structures must be provided by any IDD document claiming to implement this service.


\begin{table}[ht!]
\begin{tabularx}{\textwidth}{| p{3cm} | X |} \hline
\rowcolor{gray!33} Type & Description \\ \hline
\pdef{Address}          & A string representation of the address \\ \hline
\pdef{Boolean}          & One out of \texttt{true} or \texttt{false}. \\ \hline
\pdef{DateTime}         & Pinpoints a specific moment in time. \\ \hline
\pdef{Object}           & Set of primitives and possible further objects. \\ \hline
\pdef{Interface}        & Any suitable type chosen by the implementor of service \\ \hline
\pdef{List}$<$A$>$      & An \textit{array} of a known number of items, each having type A. \\ \hline
\pdef{Name}             & A string identifier that is intended to be both human and machine-readable. \\ \hline
\pdef{Number}           & Decimal number \\ \hline
\pdef{SecureType}       & Any suitable type chosen by the implementor of service \\ \hline
\pdef{Version}          & Specifies a service version. \\ \hline
\end{tabularx}
\end{table}

\newpage

\bibliographystyle{IEEEtran}
\bibliography{bibliography}

\newpage

\section{Revision History}
\subsection{Amendments}

\noindent\begin{tabularx}{\textwidth}{| p{1cm} | p{3cm} | p{2cm} | X | p{4cm} |} \hline
\rowcolor{gray!33} No. & Date & Version & Subject of Amendments & Author \\ \hline

1 & YYYY-MM-DD & \arrowversion & & Xxx Yyy \\ \hline
\end{tabularx}

\subsection{Quality Assurance}

\noindent\begin{tabularx}{\textwidth}{| p{1cm} | p{3cm} | p{2cm} | X |} \hline
\rowcolor{gray!33} No. & Date & Version & Approved by \\ \hline

1 & YYYY-MM-DD & \arrowversion  &  \\ \hline

\end{tabularx}

\end{document}